%%%%%%%%%%%%%%%%%%%%%%%%%%%%%%%%%%%%%%%%%%%%%%%%%%%%%%%%%%%%%%%%%%%%%%%%
%%%%%%%%%%%%%%%%%%%%%% Simple LaTeX CV Template %%%%%%%%%%%%%%%%%%%%%%%%
%%%%%%%%%%%%%%%%%%%%%%%%%%%%%%%%%%%%%%%%%%%%%%%%%%%%%%%%%%%%%%%%%%%%%%%%

%%%%%%%%%%%%%%%%%%%%%%%%%%%%%%%%%%%%%%%%%%%%%%%%%%%%%%%%%%%%%%%%%%%%%%%%
%% NOTE: If you find that it says                                     %%
%%                                                                    %%
%%                           1 of ??                                  %%
%%                                                                    %%
%% at the bottom of your first page, this means that the AUX file     %%
%% was not available when you ran LaTeX on this source. Simply RERUN  %%
%% LaTeX to get the ``??'' replaced with the number of the last page  %%
%% of the document. The AUX file will be generated on the first run   %%
%% of LaTeX and used on the second run to fill in all of the          %%
%% references.                                                        %%
%%%%%%%%%%%%%%%%%%%%%%%%%%%%%%%%%%%%%%%%%%%%%%%%%%%%%%%%%%%%%%%%%%%%%%%%

%%%%%%%%%%%%%%%%%%%%%%%%%%%% Document Setup %%%%%%%%%%%%%%%%%%%%%%%%%%%%

% Don't like 10pt? Try 11pt or 12pt
\documentclass[11pt]{article}
\RequirePackage[T1]{fontenc}
\usepackage{graphicx,amsmath,amsfonts,amssymb,amscd,yhmath,mathrsfs}
\usepackage[sorting = ydnt,defernumbers,maxbibnames=99]{biblatex}
\defbibheading{bibliography}{}
\AtBeginBibliography{\small}
\addbibresource{publications.bib}
\AtDataInput[inproceedings]{%
 \xifinlistcs{\thefield{entrykey}}{entrylist:\therefsection}{}{%
  \listcsxadd{entrylist:\therefsection}{\thefield{entrykey}}%
  \csnumgdef{inproceedingsentrycount:\therefsection}{%
      \csuse{inproceedingsentrycount:\therefsection}+1}}}

\DeclareFieldFormat[inproceedings]{labelnumber}{\mkbibinproceedingsdesc{#1}}

\newrobustcmd*{\mkbibinproceedingsdesc}[1]{%
   \number\numexpr\csuse{inproceedingsentrycount:\therefsection}+1-#1\relax}

\AtDataInput[article]{%
 \xifinlistcs{\thefield{entrykey}}{entrylist:\therefsection}{}{%
  \listcsxadd{entrylist:\therefsection}{\thefield{entrykey}}%
  \csnumgdef{articleentrycount:\therefsection}{%
      \csuse{articleentrycount:\therefsection}+1}}}

\DeclareFieldFormat[article]{labelnumber}{\mkbibarticledesc{#1}}

\newrobustcmd*{\mkbibarticledesc}[1]{%
   \number\numexpr\csuse{articleentrycount:\therefsection}+1-#1\relax}
   

\AtDataInput[misc]{%
 \xifinlistcs{\thefield{entrykey}}{entrylist:\therefsection}{}{%
  \listcsxadd{entrylist:\therefsection}{\thefield{entrykey}}%
  \csnumgdef{miscentrycount:\therefsection}{%
      \csuse{miscentrycount:\therefsection}+1}}}

\DeclareFieldFormat[misc]{labelnumber}{\mkbibmiscdesc{#1}}

\newrobustcmd*{\mkbibmiscdesc}[1]{%
   \number\numexpr\csuse{miscentrycount:\therefsection}+1-#1\relax}
   

   \AtDataInput[patent]{%
 \xifinlistcs{\thefield{entrykey}}{entrylist:\therefsection}{}{%
  \listcsxadd{entrylist:\therefsection}{\thefield{entrykey}}%
  \csnumgdef{patententrycount:\therefsection}{%
      \csuse{patententrycount:\therefsection}+1}}}

\DeclareFieldFormat[patent]{labelnumber}{\mkbibpatentdesc{#1}}

\newrobustcmd*{\mkbibpatentdesc}[1]{%
   \number\numexpr\csuse{patententrycount:\therefsection}+1-#1\relax}
   
   
      \AtDataInput[unpublished]{%
 \xifinlistcs{\thefield{entrykey}}{entrylist:\therefsection}{}{%
  \listcsxadd{entrylist:\therefsection}{\thefield{entrykey}}%
  \csnumgdef{unpublishedentrycount:\therefsection}{%
      \csuse{unpublishedentrycount:\therefsection}+1}}}

\DeclareFieldFormat[unpublished]{labelnumber}{\mkbibunpublisheddesc{#1}}

\newrobustcmd*{\mkbibunpublisheddesc}[1]{%
   \number\numexpr\csuse{unpublishedentrycount:\therefsection}+1-#1\relax}


% LaTeX will typeset using Computer Modern Roman, which a lot of
% non-mathematicians and non-engineers won't like. Also, a few PDF
% viewers may not render CMR very well. Instead, Times New Roman can
% be used. That's what this package does.
\usepackage{times}

% The automated optical recognition software used to digitize resume
% information works best with fonts that do not have serifs. This
% command uses a sans serif font throughout. Uncomment both lines (or at
% least the second) to restore a Roman font (i.e., a font with serifs).
% (NOTE: This requires the times package above)
%\renewcommand{\familydefault}{\sfdefault}

% This is a helpful package that puts math inside length specifications
\usepackage{calc}

% This package helps LaTeX auto-hyphenate hyphenated words if you use
% special hyphens. For example, bio\-/mimicry will properly hyphenate
% ``mimicry'' if necessary.
\usepackage[shortcuts]{extdash}

% Layout: Puts the section titles on left side of page
\reversemarginpar

%
%         PAPER SIZE, PAGE NUMBER, AND DOCUMENT LAYOUT NOTES:
%
% The next \usepackage line changes the layout for CV style section
% headings as marginal notes. It also sets up the paper size as either
% letter or A4. By default, letter was used. If A4 paper is desired,
% comment out the letterpaper lines and uncomment the a4paper lines.
%
% As you can see, the margin widths and section title widths can be
% easily adjusted.
%
% ALSO: Notice that the includefoot option can be commented OUT in order
% to put the PAGE NUMBER *IN* the bottom margin. This will make the
% effective text area larger.
%
% IF YOU WISH TO REMOVE THE ``of LASTPAGE'' next to each page number,
% see the note about the +LP and -LP lines below. Comment out the +LP
% and uncomment the -LP.
%
% IF YOU WISH TO REMOVE PAGE NUMBERS, be sure that the includefoot line
% is uncommented and ALSO uncomment the \pagestyle{empty} a few lines
% below.
%

%% Use these lines for letter-sized paper
% \usepackage[paper=letterpaper,
%             %includefoot, % Uncomment to put page number above margin
%             marginparwidth=1.2in,     % Length of section titles
%             marginparsep=.05in,       % Space between titles and text
%             margin=1in,               % 1 inch margins
%             includemp]{geometry}

% Use these lines for A4-sized paper
\usepackage[paper=a4paper,
           %includefoot, % Uncomment to put page number above margin
           marginparwidth=30.5mm,    % Length of section titles
           marginparsep=1.5mm,       % Space between titles and text
           margin=15mm,              % 25mm margins
           includemp]{geometry}

%% More layout: Get rid of indenting throughout entire document
\setlength{\parindent}{0in}

% Provides special list environments and macros to create new ones
\usepackage[shortlabels]{enumitem}
\usepackage{etaremune}
% Simpler bibsections for CV sections
% (thanks to natbib for inspiration)
%
% * For lists of references with hanging indents and no numbers:
%
%   \begin{bibsection}
%       \item ...
%   \end{bibsection}
%
% * For numbered lists of references (with hanging indents):
%
%   \begin{bibenum}
%       \item ...
%   \end{bibenum}
%
%   Note that bibenum numbers continuously throughout. To reset the
%   counter, use
%
%   \restartlist{bibenum}
%
%   at the place where you want the numbering to reset.

% \makeatletter
% \newlength{\bibhang}
% \setlength{\bibhang}{1em}
% \newlength{\bibsep}
%  {\@listi \global\bibsep\itemsep \global\advance\bibsep by\parsep}
% \newlist{bibsection}{itemize}{3}
% \setlist[bibsection]{label=,leftmargin=\bibhang,%
%         itemindent=-\bibhang,
%         itemsep=\bibsep,parsep=\z@,partopsep=0pt,
%         topsep=0pt}
% \newlist{bibenum}{enumerate}{3}
% \setlist[bibenum]{label=[\arabic*],resume,leftmargin={\bibhang+\widthof{[999]}},%
%         itemindent=-\bibhang,
%         itemsep=\bibsep,parsep=\z@,partopsep=0pt,
%         topsep=0pt}
% \let\oldendbibenum\endbibenum
% \def\endbibenum{\oldendbibenum\vspace{-.6\baselineskip}}
% \let\oldendbibsection\endbibsection
% \def\endbibsection{\oldendbibsection\vspace{-.6\baselineskip}}
% \makeatother

%% Reference the last page in the page number
%
% NOTE: comment the +LP line and uncomment the -LP line to have page
%       numbers without the ``of ##'' last page reference)
%
% NOTE: uncomment the \pagestyle{empty} line to get rid of all page
%       numbers (make sure includefoot is commented out above)
%
\usepackage{fancyhdr,lastpage}
\pagestyle{fancy}
%\pagestyle{empty}      % Uncomment this to get rid of page numbers
\fancyhf{}\renewcommand{\headrulewidth}{0pt}
\fancyfootoffset{\marginparsep+\marginparwidth}
\newlength{\footpageshift}
\setlength{\footpageshift}
          {0.5\textwidth+0.5\marginparsep+0.5\marginparwidth-2in}
\lfoot{\hspace{\footpageshift}%
       \parbox{4in}{\, \hfill %
                    \arabic{page} of \protect\pageref*{LastPage} % +LP
%                    \arabic{page}                               % -LP
                    \hfill \,}}

% Finally, give us PDF bookmarks
\usepackage{color,hyperref}
\definecolor{darkblue}{rgb}{0.0,0.0,0.3}
\hypersetup{colorlinks,breaklinks,
            linkcolor=blue,urlcolor=blue,
            anchorcolor=blue,citecolor=blue}

%%%%%%%%%%%%%%%%%%%%%%%% End Document Setup %%%%%%%%%%%%%%%%%%%%%%%%%%%%


%%%%%%%%%%%%%%%%%%%%%%%%%%% Helper Commands %%%%%%%%%%%%%%%%%%%%%%%%%%%%

%%% HEADING AT TOP OF CURRICULUM VITAE

% The title (name) with a horizontal rule under it
% (optional argument typesets an object right-justified across from name
%  as well)
%
% Usage: \makeheading{name}
%        OR
%        \makeheading[right_object]{name}
%
% Place at top of document. It should be the first thing.
% If ``right_object'' is provided in the square-braced optional
% argument, it will be right justified on the same line as ``name'' at
% the top of the CV. For example:
%
%       \makeheading[\emph{Curriculum vitae}]{Your Name}
%
% will put an emphasized ``Curriculum vitae'' at the top of the document
% as a title. Likewise, a picture could be included:
%
%   \makeheading[{\includegraphics[height=1.5in]{my_picture}}]{Your Name}
%
% the picture will be flush right across from the name. For this example
% to work, make sure the extra set of curly braces is included. Also
% makes ure that \usepackage{graphicx} is somewhere in the preamble.
\newcommand{\makeheading}[2][]%
        {\hspace*{-\marginparsep minus \marginparwidth}%
         \begin{minipage}[t]{\textwidth+\marginparwidth+\marginparsep}%
             {\large \bfseries #2 \hfill #1}\\[-0.15\baselineskip]%
                 \rule{\columnwidth}{1pt}%
         \end{minipage}}

%%% SECTION HEADINGS

% The section headings. Flush left in small caps down pseudo-margin.
%
% Usage: \section{section name}
\renewcommand{\section}[1]{\pagebreak[3]%
    \vspace{1.3\baselineskip}%
    \phantomsection\addcontentsline{toc}{section}{#1}%
    \noindent\llap{\scshape\smash{\parbox[t]{\marginparwidth}{\hyphenpenalty=10000\raggedright #1}}}%
    \vspace{-\baselineskip}\par}

%%% LISTS

% This macro alters a list by removing some of the space that follows the list
% (is used by lists below)
\newcommand*\fixendlist[1]{%
    \expandafter\let\csname preFixEndListend#1\expandafter\endcsname\csname end#1\endcsname
    \expandafter\def\csname end#1\endcsname{\csname preFixEndListend#1\endcsname\vspace{-0.6\baselineskip}}}

% These macros help ensure that items in outer-type lists do not get
% separated from the next line by a page break
% (they are used by lists below)
\let\originalItem\item
\newcommand*\fixouterlist[1]{%
    \expandafter\let\csname preFixOuterList#1\expandafter\endcsname\csname #1\endcsname
    \expandafter\def\csname #1\endcsname{\let\oldItem\item\def\item{\pagebreak[2]\oldItem}\csname preFixOuterList#1\endcsname}
    \expandafter\let\csname preFixOuterListend#1\expandafter\endcsname\csname end#1\endcsname
    \expandafter\def\csname end#1\endcsname{\let\item\oldItem\csname preFixOuterListend#1\endcsname}}
\newcommand*\fixinnerlist[1]{%
    \expandafter\let\csname preFixInnerList#1\expandafter\endcsname\csname #1\endcsname
    \expandafter\def\csname #1\endcsname{\let\oldItem\item\let\item\originalItem\csname preFixInnerList#1\endcsname}
    \expandafter\let\csname preFixInnerListend#1\expandafter\endcsname\csname end#1\endcsname
    \expandafter\def\csname end#1\endcsname{\csname preFixInnerListend#1\endcsname\let\item\oldItem}}

% An itemize-style list with lots of space between items
%
% Usage:
%   \begin{outerlist}
%       \item ...    % (or \item[] for no bullet)
%   \end{outerlist}
\newlist{outerlist}{itemize}{3}
    \setlist[outerlist]{label=\enskip\textbullet,leftmargin=*}
    \fixendlist{outerlist}
    \fixouterlist{outerlist}

% An environment IDENTICAL to outerlist that has better pre-list spacing
% when used as the first thing in a \section
%
% Usage:
%   \begin{lonelist}
%       \item ...    % (or \item[] for no bullet)
%   \end{lonelist}
\newlist{lonelist}{itemize}{3}
    \setlist[lonelist]{label=\enskip\textbullet,leftmargin=*,partopsep=0pt,topsep=0pt}
    \fixendlist{lonelist}
    \fixouterlist{lonelist}

% An itemize-style list with little space between items
%
% Usage:
%   \begin{innerlist}
%       \item ...    % (or \item[] for no bullet)
%   \end{innerlist}
\newlist{innerlist}{itemize}{3}
    \setlist[innerlist]{label=\enskip\textbullet,leftmargin=*,parsep=0pt,itemsep=0pt,topsep=0pt,partopsep=0pt}
    \fixinnerlist{innerlist}

% An environment IDENTICAL to innerlist that has better pre-list spacing
% when used as the first thing in a \section
%
% Usage:
%   \begin{loneinnerlist}
%       \item ...    % (or \item[] for no bullet)
%   \end{loneinnerlist}
\newlist{loneinnerlist}{itemize}{3}
    \setlist[loneinnerlist]{label=\enskip\textbullet,leftmargin=*,parsep=0pt,itemsep=0pt,topsep=0pt,partopsep=0pt}
    \fixendlist{loneinnerlist}
    \fixinnerlist{loneinnerlist}

%%% EXTRA SPACE

% To add some paragraph space between lines.
% This also tells LaTeX to preferably break a page on one of these gaps
% if there is a needed pagebreak nearby.
\newcommand{\blankline}{\quad\pagebreak[3]}
\newcommand{\halfblankline}{\quad\vspace{-0.5\baselineskip}\pagebreak[3]}

%%% FORMATTING MACROS

% Provides a linked \doi{#1} that links doi:#1 to http://dx.doi.org/#1
\usepackage{doi}
% To change the text before the DOI, adjust this command
%\renewcommand\doitext{doi:}

% Provides a linked \url{#1} that doesn't require escape characters
\usepackage{url}

% You can adjust the style \url{} uses here:
% (options are: same, rm, sf, tt; defaults to tt)
\urlstyle{same}

% For \email{ADDRESS}, links ADDRESS to the url mailto:ADDRESS
% (uncomment to typeset the e\-/mail address in typewriter font;
%  otherwise, will be typeset in the \urlstyle above)
%\DeclareUrlCommand\emaillink{\urlstyle{tt}}
\providecommand*\emaillink[1]{\nolinkurl{#1}}
\providecommand*\email[1]{\href{mailto:#1}{\emaillink{#1}}}

\providecommand\BibTeX{{B\kern-.05em{\sc i\kern-.025em b}\kern-.08em \TeX}}
\providecommand\Matlab{\textsc{Matlab}}
\newcommand*{\nolink}[1]{%
  {\protect\NoHyper#1\protect\endNoHyper}%
}
% Custom hyphenation rules for words that LaTeX has trouble with
% \hyphenation{bio-mim-ic-ry bio-in-spi-ra-tion re-us-a-ble pro-vid-er Media-Wiki}

%%%%%%%%%%%%%%%%%%%%%%%% End Helper Commands %%%%%%%%%%%%%%%%%%%%%%%%%%%

%%%%%%%%%%%%%%%%%%%%%%%%% Begin CV Document %%%%%%%%%%%%%%%%%%%%%%%%%%%%

\begin{document}

\begin{refsection}
\makeheading{\huge Filippo Pecci}

\section{Contact Information}

% NOTE: Mind where the & separators and \\ breaks are in the following
%       table. Table is one row made up of three parboxes. The left
%       parbox has address info, the middle parbox has a vertical bar,
%       and the right parbox has phone and electronic contact
%       information.
%
% MACROS: \rcollength is the width of the right column of the table
%             (adjust it to your liking; default is 1.85in).
%         \spacewidth is width of area between left and right boxes.
%
\newlength{\rcollength}\setlength{\rcollength}{1.85in}%
\newlength{\spacewidth}\setlength{\spacewidth}{20pt}
%
\begin{tabular}[t]{@{}p{\textwidth-\rcollength-\spacewidth}@{}p{\spacewidth}@{}p{\rcollength}}%

% Address box
\parbox{\textwidth-\rcollength-\spacewidth}{%
%Assistant Professor\\
{Princeton University}\\
{Andlinger Center for Energy and the Environment}\\
Princeton, NJ 08540 USA}

&
% Uncomment to add a vertical bar in middle of contact information
%{\vrule width 0.5pt}
\parbox[m][5\baselineskip]{\spacewidth}{} &

% Non-snail-mail contact information
\parbox{\rcollength}{%
%\textit{Work:} +1-480-965-2899 \\
%\textit{Fax:} +1-480-965-2751 \\
\email{filippopecci@princeton.edu}\\
\href{https://filippopecci.github.io/}{Academic website}}

\end{tabular}

%%
%% In modern CV's, it seems like ``Objective'' is frowned upon. Instead,
%% incorporate it into a well-constructed cover letter. The ``More
%% information'' can go at the end of the CV, but it should not distract
%% from the section giving references available to contact.
%%
%
% \section{Objective}
%
% Placement in an academic position (i.e., faculty, postdoctoral, or
% research scientist) that allows for advanced research in distributed
% complex adaptive systems (i.e., modeling, analysis, design, and
% verification) with a particular focus on the control of engineered
% agents (e.g., for communications, control, software, electronics, and
% sustainability) and the analysis of biological phenomena (e.g.,
% self-organization, ecological rationality)
% \begin{innerlist}
% \item More information and auxiliary documents can be found at\\\url{http://www.tedpavlic.com/facjobsearch/}
% \end{innerlist}

\section{Research objective}
Develop and apply computational methods to solve large-scale mixed-integer optimization models, optimize design and control of complex network systems, and provide decision support to accelerate transition to net-zero emissions and resilient infrastructure systems.

\section{Appointments}
\begin{tabular}{p{0.15\linewidth}  p{0.75\linewidth}}
    2022-present & Associate Research Scholar, Andlinger Center for Energy \& the Environment, Princeton University \\
    \hfill\\
    2018-2022 & Postdoctoral Research Associate, Department of Civil \& Environmental Engineering, Imperial College London
\end{tabular}

\halfblankline

%\halfblankline

% \begin{innerlist}

%     \item[] \href{http://sols.asu.edu/}{School of Life Sciences}
%     \begin{innerlist}
%         \item Supervisor: \href{http://www.public.asu.edu/~spratt1}{Professor Stephen C.~Pratt}
%         \item Decentralized decision making and behavioral bio-mimicry
%             of social insects
%     \end{innerlist}

% \end{innerlist}

\section{Education}
\begin{tabular}{p{0.15\linewidth}  p{0.75\linewidth}}
    2018 & Ph.D. in Computational Optimization, Department of Civil \& Environmental Engineering, Imperial College London \\
    & Thesis: Optimal design for control of water supply networks by mixed integer programming [\href{http://hdl.handle.net/10044/1/72854}{pdf}]\\
    \halfblankline \\
    2014 & MSc in Mathematics, Università degli Studi di Padova (Italy)\\
    %& Thesis: Hierarchical stratification of Pareto sets [\href{https://arxiv.org/abs/1407.1755}{pdf}]\\
    \halfblankline \\
    2011 & BSc in Mathematics, Università degli Studi di Padova (Italy)
\end{tabular}

% \section{Submitted Journal Publications}
%
% % Add a little space to nudge next ``Ref'd Journal Publications'' marginpar
% % down to make room for tall ``Submitted Journal Publications''
% % marginpar. If there are enough submitted journal publications, this
% % space will not be needed (and should be removed).
% \vspace{0.1in}

\section{Teaching \& Supervision}
Princeton University
    \begin{innerlist}
        \item Guest lecturer for Applied Optimization Methods for Energy Systems Engineering, Fall 2022 (500 level graduate elective). Lecture on decomposition methods with application to energy systems.\\
    \end{innerlist}
{Imperial College London}
\begin{innerlist}
     \item Main lecturer for BSc and MSc modules on convex optimization with applications in water systems, Fall \& Spring 2020 and 2021.
     \item Guest lecturer and teaching assistant for MSc modules on modelling and optimization of water distribution networks, Spring 2016, 2017, 2018, 2019, 2020, and 2021.
     \item Mathematics tutor for MSc students, with focus on linear algebra and calculus, Fall 2015, 2016, 2017.
     \item Technical supervision of 4 PhD students working on optimal design, model identification and fault estimation problems for water distribution networks.
    \item Co-supervision of 4 final year MSc projects on analysis and optimization of water distribution networks.
\end{innerlist}

\nocite{*}
\section{Preprints}
\newrefcontext[labelprefix=P]
\printbibliography[type = unpublished]

\section{Refeered journal articles}
\newrefcontext[labelprefix=J]
\printbibliography[type =article]

\section{{Refereed conference proceedings}}
\newrefcontext[labelprefix=C]
\printbibliography[type = inproceedings]

\section{{Patents}}
\newrefcontext[labelprefix=B]
\printbibliography[type = patent]
%\section{{Others}}
%\newrefcontext[labelprefix=R]
%\printbibliography[type = review]
% \newrefcontext[labelprefix=A]
% \printbibliography[type = misc]
\section{Conference presentations and invited seminars}
\small
\begin{etaremune}
\item 2023 INFORMS Annual Meeting, Phoenix (Arizona), 15-18 October, 2023. Learning to optimize macro-energy systems.
\item  International Conference on Optimization and Decision Science 2022, Florence, Italy, 30 August - 2 September, 2022. A global optimization framework for resilient water distribution networks. [\href{https://filippopecci.github.io/uploads/slides/ODS_2022.pdf}{slides}]
\item European Control Conference 2022, London, United Kingdom, 12-15 July, 2022. Optimal Design-for-Control of Chlorine Booster Systems in Water Networks via Convex Optimization. [\href{https://filippopecci.github.io/uploads/slides/ODS_2022.pdf}{slides}]
\item Control \& optimization Seminars, Imperial College London, 22 Gennaio 2020.  Mathematical optimization for intelligent water distribution networks: model calibration, and event detection and localisation.
\item 17th Computing and Control for the Water Industry (CCWI), Exeter, United Kingdom, 1-4 September, 2019. Tight Convex Relaxations for Optimal Design and Control Problems in Water distribution Networks. [\href{https://filippopecci.github.io/uploads/slides/CCWI_2019.pdf}{slides}] 
\item 6th International Conference on Continuous Optimization (ICCOPT), Berlin, Germany, 3 - 8 August, 2019.  Non-linear inverse problems via sequential convex optimization. [\href{https://filippopecci.github.io/uploads/slides/ICCOPT_2019.pdf}{slides}]
\item  6th International Conference on Engineering Optimization (EngOpt), Lisbon, Portugal, 17 - 19 September, 2018. A branch and bound method for globally optimizing valve locations in water distribution networks.[\href{https://filippopecci.github.io/uploads/slides/ENGOPT_2018.pdf}{slides}]
\item  20th IFAC World Congress, Toulouse, France, 9 - 14 July, 2017. Outer approximation methods for the solution of co-design optimization problems in water distribution networks. [\href{https://filippopecci.github.io/uploads/slides/IFAC_2017.pdf}{slides}]
\item 14th Computing and Control for the Water Industry (CCWI), Amsterdam, the Netherlands, 7-9 November, 2016. Multiobjective pressure optimization in water distribution systems (Poster Presentation).
\item 13th Computing and Control for the Water Industry (CCWI), Leicester, United Kingdom, 2-4 September, 2015. Mathematical programming methods for pressure management in water distribution systems. [\href{https://filippopecci.github.io/uploads/slides/CCWI_2015.pdf}{slides}]
\end{etaremune}
\end{refsection}


\end{document}
